\chapter{
مبانی زبان 
\eng{Haskell}
}
\label{chp:chap2}\thispagestyle{empty}
%==========+==========+==========+==========+==========+==========+==========+==========
در این فصل به آموزش مقدماتی زبان 
\eng{Haskell}
میپردازیم. در طول این فصل مباحثی از 
\eng{Haskell}
را بیان میکنیم که پیش نیازی از فصل های بعدی اند. این فصل دو قسمت است. در قسمت اول به آموزش زبان به صورت نظری میپردازیم و در بخش دوم از آموزه های بخش اول استفاده میکنیم تا تعدادی برنامه ساده بنویسیم. هر برنامه ای که در آن قسمت نوشته شود میتواند توسط خواننده نوشته و اجرا شود. 
\section{
مقدماتی از استاندارد زبان
}
یک شناسه یا 
\eng{identifier}
ار یک حرف و بعد از آن هیچ یا چند حرف، زقم، 
\eng{\_}
، علامت نقل قول مفرد 
\eng{\'}
باید ساخته شود. شناسه ها به صورت لغوی به دو فضای نام 
\eng{namespace}
تقسیم میشوند: آنهایی که حرف اول آنها کوچک برای متغیر ها و آنهایی که حرف اول آنها بزرگ است برای سازنده ها
\eng{constructor}
 هستند. بزای مثال سازنده 
\hs{Int}
برای ساخت عداد صحیح با حرف اول بزرگ مشخص میشود. نام گذاری باقی بخش های برنامه قوانین از پیش تایین شده ای دارند. این قوانین به تفصیل در
\href{https://www.haskell.org/onlinereport/haskell2010/}{\eng{Haskell Report 2010}}
بحث شده اند. لازم به ذکر است که کامپایلر
\eng{GHC}
که برای این گزارش استفاده شده، دارای توسیع هایی از این استاندارد است که تغییراتی جزیی در آن ایجاد کرده اند. تفاوت ها را میتوان در 
\href{https://www.haskell.org/ghc/}{\eng{haskell.org}}
در صفحه شناسه مربوط به هر نسخه کامپایلر دید.
\newline
نام گذاری در زبان 
\eng{Haskell}
به بزرگی و کوچکی حروف حساس است. به عنوان مثال
\hs{name}
، 
\hs{naMe}
و
\hs{Name}
همگی شناسه های متفاوتی هستند. به طور خلاصه اسامی توابع و متغیر ها باید با حروف کوچک شروع شوند اما بعد از آن برنامه نویس در انتخاب بقیه شناسه آزاد است.
\newline
در زبان 
\eng{Haskell}
میتوان از نوشتن نقطه ویرگول 
\hs{;}
و آکولاد
\hs{\{ \}}
صرف نظر و از تورفتگی یا 
\eng{indentation}
برای چیدمان برنامه استفاده کرد. وقتی از چیدمان استفاده شود باید قوانینی رعایت شوند؛ در غیر این صورت کامپایلر با برگرداندن خطا متوقف میشود. به صورت غیر رسمی در جایی که از آکولاد و نقطه ویرگول صرف نظر شده باشد زبان به صورت زیر عمل میکند: